\documentclass[10pt,a4paper]{article}
\usepackage{amsmath,amssymb,amsthm,graphicx}
\usepackage{amsmath,amssymb,amsthm,mathtools,hyperref,cleveref}
\newtheorem{theorem}{Theorem}
\newtheorem{lemma}[theorem]{Lemma}
\newtheorem{proposition}[theorem]{Proposition}
\newtheorem{corollary}[theorem]{Corollary}
\newtheorem{definition}[theorem]{Definition}
\newtheorem{conjecture}[theorem]{Conjecture}
\newtheorem{remark}[theorem]{Remark}
\newcommand{\lean}[1]{}
\newcommand{\proves}[1]{}
\newcommand{\uses}[1]{}
\newcommand{\leanok}{}
\newcommand{\chapter}[1]{\section{#1}}
% Put any macro and import needed for the project here.
% This will be used by both the web and print versions of the blueprint.
% This file is not meant to be built. Build src/web.tex or src/print.text instead.

% Letters
\newcommand{\C}{\mathbb{C}}
\newcommand{\Z}{\mathbb{Z}}
\newcommand{\bbc}{\mathbb{C}}
\newcommand{\E}{\mathbb{E}}
\newcommand*{\bbe}{\mathbb{E}}
\newcommand{\F}{\mathbb{F}}
\newcommand{\bbf}{\mathbb{F}}
\newcommand{\bbH}{\mathbb{H}}
\newcommand{\bbP}{\mathbb{P}}
\newcommand{\bbI}{\mathbb{I}}
\newcommand{\bbn}{\mathbb{N}}
\newcommand{\bbq}{\mathbb{Q}}
\newcommand{\bbr}{\mathbb{R}}
\newcommand{\bbt}{\mathbb{T}}
\newcommand{\bbz}{\mathbb{Z}}
\newcommand{\N}{\mathbb{N}}

\newcommand{\lo}[1]{\mathcal{L}{#1}}

% Paired delimiters
\newcommand{\abs}[1]{\left\lvert #1\right\rvert}
\newcommand{\Abs}[1]{\lvert #1 \rvert}
\newcommand{\brac}[1]{\left( #1\right)}
\newcommand{\norm}[1]{\lVert #1\rVert}
\newcommand{\inn}[1]{\left\langle #1 \right\rangle}

% Operators
\DeclareMathOperator{\dist}{dist}

\newcommand{\ind}[1]{1_{#1}}
\providecommand{\tup}[1]{{\vec{#1}}}
\title{Asymmetric weak Polynomial Freiman-Ruzsa over the integers}
\author{Terence Tao, PFR project contributors}
\date{github:teorth-pfr:4c56de547b2f33b0bc47b92f3bfbcebc6457a776; Apache-2.0}
\begin{document}
\maketitle
\noindent\textit{Editorial note.} Source-organized author material only. All mathematical source spans are retained unchanged. Blueprint lean/proves/uses metadata is hidden as formatting only; numbering is local to this material wrapper.\par\medskip
\noindent\textit{Editorial note.} Ordinary cited prerequisites and concise author arguments are accepted as such. No independent mathematical certification or public ID allocation is claimed.\par\medskip
\noindent\textit{Allowed standard prerequisites (editorial).}\begin{itemize}
\item Standard finite Shannon entropy identities, conditional entropy and finite abelian group algebra are cited in their original Blueprint formulations.
\item Named prior PFR/entropy or BSG results are allowed prerequisites; their original statements are retained where used, without demanding recursively expanded proofs.
\end{itemize}\medskip
\chapter{The Fibring lemma}

\begin{proposition}[General fibring identity]\label{fibring-ident}
  \lean{rdist_of_indep_eq_sum_fibre, rdist_le_sum_fibre}\leanok
  Let $\pi : H \to H'$ be a homomorphism additive groups, and let $Z_1,Z_2$ be $H$-valued random variables. Then we have
  \[
    d[Z_1; Z_2] \geq d[\pi(Z_1);\pi(Z_2)] + d[Z_1|\pi(Z_1); Z_2 |\pi(Z_2)].
  \]
  Moreover, if $Z_1,Z_2$ are taken to be independent, then the difference between the two sides is
$$I( Z_1 - Z_2 : (\pi(Z_1), \pi(Z_2))  |  \pi(Z_1 - Z_2) ).$$
\end{proposition}


\begin{proof}\uses{ruz-copy, independent-exist, submodularity,conditional-mutual-alt,chain-rule,relabeled-entropy, cond-dist-alt}\leanok
  Let $Z_1,Z_2$ be independent throughout (this is possible by \Cref{ruz-copy} and \Cref{independent-exist}).  By \Cref{cond-dist-alt}, We have
  \begin{align*}
    & d[Z_1  |\pi(Z_1); Z_2 |\pi(Z_2)] \\
    & = \bbH[Z_1 - Z_2 | \pi(Z_1),\pi(Z_2)] - \tfrac{1}{2} \bbH[Z_1 | \pi(Z_1)] - \tfrac{1}{2} \bbH[Z_2 | \pi(Z_2)] \\
    & \leq  \bbH[Z_1 - Z_2 | \pi(Z_1+Z_2)]  - \tfrac{1}{2} \bbH[Z_1 | \pi(Z_1)] - \tfrac{1}{2}H[Z_2 | \pi(Z_2)] \\
    & = d[Z_1;Z_2] - d[\pi(Z_1);\pi(Z_2)].
  \end{align*}
  In the middle step, we used \Cref{submodularity}, and in the last step we used the fact that
  \[\bbH[Z_1 - Z_2 |  \pi(Z_1-Z_2)] = \bbH[Z_1 - Z_2] - \bbH[\pi(Z_1-Z_2)]\]
  (thanks to \Cref{chain-rule} and \Cref{relabeled-entropy}) and that
  \[\bbH[Z_i| \pi(Z_i)] = \bbH[Z_i] - \bbH[\pi(Z_i)]\] (since $Z_i$ determines $\pi(Z_i)$).
  This gives the claimed inequality. The difference between the two sides is precisely
  \[\bbH[Z_1 - Z_2  | \pi(Z_1 - Z_2)] - \bbH[Z_1 - Z_2  | \pi(Z_1),\pi(Z_2)].\]
  To rewrite this in terms of (conditional) mutual information, we use the identity
  \[\bbH[A|B] - \bbH[A | B,C] = \bbI[A : C | B],\]
  (which follows \Cref{conditional-mutual-alt})
  taking
  $A := Z_1 - Z_2$, $B := \pi(Z_1 - Z_2)$ and $C := (\pi(Z_1),\pi(Z_{2}))$, and noting that in this case $\bbH[A | B,C] = \bbH[A | C]$ since $C$ uniquely determines $B$ (this may require another helper lemma about entropy).
  This completes the proof.
\end{proof}


\begin{corollary}\label{fibring-ineq}
 \lean{rdist_of_hom_le} \leanok
If $\pi:G\to H$ is a homomorphism of additive groups and $X,Y$ are $G$-valued random variables then
\[d[X;Y]\geq d[\pi(X);\pi(Y)].\]
\end{corollary}

\begin{proof}
\uses{fibring-ident, ruzsa-nonneg}\leanok
By \Cref{fibring-ident} and the nonnegativity of conditional Ruzsa distance (from \Cref{ruzsa-nonneg}) we have
\[d[X;Y]\geq d[\pi(X);\pi(Y)]+d[X\mid \pi(X);Y\mid \pi(Y)].\]
The inequality follows from $d[X\mid \pi(X);Y\mid \pi(Y)]\geq 0$ (\Cref{ruzsa-nonneg}).
\end{proof}

\chapter{Shannon entropy inequalities}

Random variables in this paper are measurable maps $X : \Omega \to S$ from a probability space $\Omega$ to a measurable space $S$, and called $S$-valued random variables. In many cases we will assume that singletons in $S$ are measurable.  Often we will restrict further to the case when $S$ is finite with the discrete $\sigma$-algebra, which of course implies that $S$ has measurable singletons.

\begin{definition}[Entropy]
  \label{entropy-def}
  \lean{ProbabilityTheory.entropy}
  \leanok
  If $X$ is an $S$-valued random variable, the entropy $\bbH[X]$ of $X$ is defined
  $$ \bbH[X] := \sum_{s \in S} \bbP[X=x] \log \frac{1}{\bbP[X=x]}$$
  with the convention that $0 \log \frac{1}{0} = 0$.
\end{definition}

\begin{lemma}[Entropy and relabeling]\label{relabeled-entropy}\lean{ProbabilityTheory.entropy_comp_of_injective, ProbabilityTheory.entropy_of_comp_eq_of_comp} \uses{entropy-def}\leanok
  \begin{itemize}
\item[(i)]   If $X: \Omega \to S$ and $Y: \Omega \to T$ are random variables, and $Y = f(X)$ for some injection $f: S \to T$, then $\bbH[X] = \bbH[Y]$.
\item[(ii)]   If $X: \Omega \to S$ and $Y: \Omega \to T$ are random variables, and $Y = f(X)$ and $X = g(Y)$ for some functions $f: S \to T$, $g: T \to S$, then $\bbH[X] = \bbH[Y]$.
\end{itemize}
\end{lemma}

\begin{proof} \leanok Expand out both entropies and rearrange.
\end{proof}

\begin{lemma}[Jensen bound]\label{jensen-bound}
  \uses{entropy-def}
  \lean{ProbabilityTheory.entropy_le_log_card, ProbabilityTheory.entropy_le_log_card_of_mem}
  \leanok
  If $X$ is an $S$-valued random variable, then $\bbH[X] \leq \log |S|$.
\end{lemma}

\begin{proof}
  \uses{concave}
  \leanok

  This is a direct consequence of \Cref{concave} and Jensen's inequality.
\end{proof}


We use $X,Y$ to denote the pair $\omega \mapsto (X(\omega),Y(\omega)$).

\begin{lemma}[Commutativity and associativity of joint entropy]
  \label{entropy-comm}
  \lean{ProbabilityTheory.entropy_comm, ProbabilityTheory.entropy_assoc}
  \leanok
  If $X: \Omega \to S$, $Y: \Omega \to T$, and $Z: \Omega \to U$ are random variables, then $\bbH[X, Y] = \bbH[Y, X]$ and $\bbH[X, (Y,Z)] = \bbH[(X,Y), Z]$.
\end{lemma}
\begin{proof}
  \leanok
  \uses{relabeled-entropy}
  Set up an injection from $(X,Y)$ to $(Y,X)$ and use \Cref{relabeled-entropy} for the first claim. Similarly for the second claim.
\end{proof}


Note: it may happen that $E$ has zero measure.  In which case, the ambient probability measure should be replaced with a zero measure.  (In our formalization we achieve this by working with arbitrary measures, and normalizing them to be probability measures if possible, else using the zero measure.  Conditioning is also formalized using existing Mathlib definitions.)

\begin{definition}[Conditional entropy]
  \label{conditional-entropy-def}
  \uses{entropy-def,condition-event-def}
  \lean{ProbabilityTheory.condEntropy}
  \leanok
  If $X: \Omega \to S$ and $Y: \Omega \to T$ are random variables, the conditional entropy $\bbH[X|Y]$ is defined as
  $$ \bbH[X|Y] := \sum_{y \in Y} \bbP[Y = y] \bbH[(X | Y=y)].$$
\end{definition}

\begin{lemma}[Conditional entropy and relabeling]\label{relabeled-entropy-cond}
  \uses{conditional-entropy-def}
  \lean{ProbabilityTheory.condEntropy_comp_of_injective, ProbabilityTheory.condEntropy_of_injective'}\leanok
  If $X: \Omega \to S$, $Y: \Omega \to T$, and $Z: \Omega \to U$ are random variables, and $Y = f(X,Z)$ almost surely for some map $f: S \times U \to T$ that is injective for each fixed $U$, then $\bbH[X|Z] = \bbH[Y|Z]$.

  Similarly, if $g: T \to U$ is injective, then $\bbH[X|g(Y)] = \bbH[X|Y]$.
\end{lemma}

\begin{proof}
  \uses{relabeled-entropy}\leanok
  For the first part, use \Cref{conditional-entropy-def} and then \Cref{relabeled-entropy}.  The second part is a direct computation.
\end{proof}


\begin{lemma}[Chain rule]\label{chain-rule}
  \uses{conditional-entropy-def}
  \lean{ProbabilityTheory.chain_rule, ProbabilityTheory.chain_rule'}
  \leanok
  If $X: \Omega \to S$ and $Y: \Omega \to T$ are random variables, then
  $$ \bbH[X, Y] = \bbH[Y] + \bbH[X|Y].$$
  \end{lemma}
  \begin{proof}
  \leanok
  Direct computation.
\end{proof}


\begin{lemma}[Alternative formulae for mutual information]
  \label{alternative-mutual}
  \uses{information-def}
  \lean{ProbabilityTheory.mutualInfo_comm, ProbabilityTheory.mutualInfo_eq_entropy_sub_condEntropy}
  \leanok
  With notation as above, we have
  $$  \bbI[X : Y] = \bbI[Y:X]$$
  $$  \bbI[X : Y] = \bbH[X] - \bbH[X|Y]$$
  $$  \bbI[X : Y] = \bbH[Y] - \bbH[Y|X]$$
\end{lemma}

\begin{proof}
  \uses{entropy-comm,chain-rule}
  \leanok
  Immediate from Lemmas \ref{entropy-comm}, \ref{chain-rule}.
\end{proof}

\begin{lemma}[Nonnegativity of mutual information]
  \label{mutual-nonneg}
  \lean{ProbabilityTheory.mutualInfo_nonneg}
  \leanok
  We have $\bbI[X:Y] \geq 0$.
\end{lemma}

\begin{proof}
  \uses{concave, alternative-mutual}
  \leanok

  An application of jensen's inequality and \Cref{concave,alternative-mutual}.
\end{proof}

\begin{corollary}[Subadditivity]
  \label{subadditive}
  \lean{ProbabilityTheory.entropy_pair_le_add}
  \leanok
  With notation as above, we have $\bbH[X,Y] \leq \bbH[X] + \bbH[Y]$.
\end{corollary}

\begin{proof}
  \uses{mutual-nonneg, alternative-mutual}
  \leanok
  Use \Cref{mutual-nonneg}.
\end{proof}

\begin{corollary}[Conditioning reduces entropy]
  \label{cond-reduce}
  \lean{ProbabilityTheory.condEntropy_le_entropy}
  \leanok
  With notation as above, we have $\bbH[X|Y] \leq \bbH[X]$.
\end{corollary}
\begin{proof}
  \uses{mutual-nonneg, alternative-mutual}
  \leanok
  Combine \Cref{mutual-nonneg} with \Cref{alternative-mutual}.
\end{proof}

\begin{corollary}[Submodularity]\label{submodularity}
  \lean{ProbabilityTheory.entropy_submodular, ProbabilityTheory.condEntropy_comp_ge}\leanok
  With three random variables $X,Y,Z$, one has $\bbH[X|Y,Z] \leq \bbH[X|Z]$.
\end{corollary}

\begin{proof} \uses{cond-reduce}\leanok Apply the ``averaging over conditioning'' argument to \Cref{cond-reduce}.
\end{proof}

\begin{corollary}[Alternate form of submodularity]\label{alt-submodularity}
  \lean{ProbabilityTheory.entropy_triple_add_entropy_le}\leanok
  With three random variables $X,Y,Z$, one has
  $$ \bbH[X,Y,Z] + \bbH[Z] \leq \bbH[X,Z] + \bbH[Y,Z].$$
\end{corollary}

\begin{proof}  \uses{submodularity,chain-rule} \leanok Apply \Cref{submodularity} and \Cref{chain-rule}.
\end{proof}


\begin{lemma}[Vanishing of mutual information]
  \label{vanish-entropy}
  \uses{information-def, independent-def}
  \lean{ProbabilityTheory.mutualInfo_eq_zero}
  \leanok
  If $X,Y$ are random variables, then $\bbI[X:Y] = 0$ if and only if $X,Y$ are independent.
  \end{lemma}

  \begin{proof}
    \uses{concave}
    \leanok

    An application of the equality case of Jensen's inequality and \Cref{concave}.
  \end{proof}

\begin{corollary}[Additivity of entropy]\label{add-entropy}
  \lean{ProbabilityTheory.entropy_pair_eq_add}
  \leanok
  If $X,Y$ are random variables, then $\bbH[X,Y] = \bbH[X] + \bbH[Y]$ if and only if $X,Y$ are independent.
\end{corollary}

\begin{proof} \uses{vanish-entropy}\leanok
  Direct from \Cref{vanish-entropy}.
\end{proof}


\begin{lemma}[Alternate formula for conditional mutual information]
  \label{conditional-mutual-alt}
  \uses{conditional-mutual-def}
  \lean{ProbabilityTheory.condMutualInfo_eq, ProbabilityTheory.condMutualInfo_eq'}\leanok
We have
  $$ \bbI[X:Y|Z] := \bbH[X|Z] + \bbH[Y|Z] - \bbH[X,Y|Z]$$
and
  $$ \bbI[X:Y|Z] := \bbH[X|Z] - \bbH[X|Y,Z].$$
\end{lemma}

\begin{proof}\leanok Routine computation.
\end{proof}

\chapter{Applications of Jensen's inequality}

In this chapter, $h$ denotes the function $h(x) := x \log \frac{1}{x}$ for $x \in [0,1]$.

\begin{lemma}[Concavity]\label{concave}
  \lean{Real.strictConcaveOn_negMulLog}\leanok
  $h$ is strictly concave on $[0,\infty)$.
\end{lemma}

\begin{proof} \leanok Check that $h'$ is strictly monotone decreasing.
\end{proof}

\chapter{Entropic Ruzsa calculus}

In this section $G$ will be a finite additive group.  (May eventually want to generalize to infinite $G$.)

\begin{lemma}[Negation preserves entropy]\label{neg-ent}
  \lean{ProbabilityTheory.entropy_neg}\leanok
  If $X$ is $G$-valued, then $\bbH[-X]=\bbH[X]$.
\end{lemma}

\begin{proof}
  \uses{relabeled-entropy}\leanok Immediate from \Cref{relabeled-entropy}.
\end{proof}

\begin{lemma}[Shearing preserves entropy]\label{shear-ent}
  \lean{ProbabilityTheory.condEntropy_add_right, ProbabilityTheory.condEntropy_add_left, ProbabilityTheory.condEntropy_sub_right, ProbabilityTheory.condEntropy_sub_left}\leanok
  If $X,Y$ are $G$-valued, then $\bbH[X \pm Y | Y]=\bbH[X|Y]$ and $\bbH[X \pm Y, Y] = \bbH[X, Y]$.
\end{lemma}

\begin{proof}
  \uses{relabeled-entropy-cond,chain-rule}\leanok Immediate from \Cref{relabeled-entropy-cond} and \Cref{chain-rule}.
\end{proof}

\begin{lemma}[Lower bound of sumset]\label{sumset-lower-gen}
  \lean{ProbabilityTheory.max_entropy_sub_mutualInfo_le_entropy_add, ProbabilityTheory.max_entropy_sub_mutualInfo_le_entropy_sub}\leanok
  If $X,Y$ are $G$-valued random variables on $\Omega$, we have
$$ \max(\bbH[X], \bbH[Y]) -  \bbI[X:Y] \leq \bbH[X \pm Y].$$
\end{lemma}

\begin{proof}\uses{cond-reduce, alternative-mutual, shear-ent, neg-ent} \leanok
    By \Cref{cond-reduce}, \ref{shear-ent}, \ref{alternative-mutual}, \ref{neg-ent} we have
$$
 \bbH[X\pm Y] \geq \bbH[X\pm Y|Y] = \bbH[X|Y]= \bbH[X] -  \bbI[X:Y]
$$
and similarly with the roles of $X,Y$ reversed, giving the claim.
\end{proof}


\begin{corollary}[Independent lower bound on sumset]
  \label{sumset-lower}
  \lean{ProbabilityTheory.max_entropy_le_entropy_add, ProbabilityTheory.max_entropy_le_entropy_sub}
  \leanok

  If $X,Y$ are independent $G$-valued random variables, then
  $$\max(\bbH[X], \bbH[Y]) \leq \bbH[X\pm Y]. $$
\end{corollary}
\begin{proof}
  \uses{sumset-lower-gen, vanish-entropy}
  \leanok

  Combine \Cref{sumset-lower-gen} with \Cref{vanish-entropy}.
\end{proof}

One random variable is said to be a copy of another if they have the same distribution.

\begin{lemma}[Copy preserves entropy]\label{copy-ent}
  \uses{entropy-def}
  \lean{ProbabilityTheory.IdentDistrib.entropy_congr}\leanok
  If $X'$ is a copy of $X$ then $\bbH[X'] = \bbH[X]$.
\end{lemma}

\begin{proof}\leanok Immediate from \Cref{entropy-def}.
\end{proof}

\begin{lemma}[Existence of independent copies]\label{independent-exist}
  \uses{independent-def}
  \lean{ProbabilityTheory.independent_copies, ProbabilityTheory.independent_copies', ProbabilityTheory.independent_copies_two, ProbabilityTheory.independent_copies3_nondep} \leanok
  Let $X_i : \Omega_i \to S_i$ be random variables for $i=1,\dots,k$.  Then if one gives $\prod_{i=1}^k S_i$ the product measure of the laws of $X_i$, the coordinate functions $(x_j)_{j=1}^k \mapsto x_i$ are jointly independent random variables which are copies of the $X_1,\dots,X_k$.
\end{lemma}

\begin{proof} \leanok
  Explicit computation.
\end{proof}

\begin{definition}[Ruzsa distance]\label{ruz-dist-def}
  \uses{entropy-def, independent-exist}
  \lean{rdist_def}\leanok
  Let $X,Y$ be $G$-valued random variables (not necessarily on the same sample space).  The \emph{Ruzsa distance} $d[X ;Y]$ between $X$ and $Y$ is defined to be
$$ d[X ;Y] := \bbH[X' - Y'] - \bbH[X']/2 - \bbH[Y']/2$$
where $X',Y'$ are (the canonical) independent copies of $X,Y$ from \Cref{independent-exist}.
\end{definition}


\begin{lemma}[Distance from zero]\label{dist-zero}
  \uses{ruz-dist-def}\lean{rdist_zero_eq_half_ent}\leanok
If $X$ is a $G$-valued random variable and $0$ is the random variable taking the value $0$ everywhere then
\[d[X;0]=\mathbb{H}(X)/2.\]
\end{lemma}

\begin{proof}\leanok
This is an immediate consequence of the definitions and $X-0\equiv X$ and $\mathbb{H}(0)=0$.
\end{proof}

\begin{lemma}[Copy preserves Ruzsa distance]\label{ruz-copy}
  \uses{ruz-dist-def}
  \lean{ProbabilityTheory.IdentDistrib.rdist_congr}\leanok
  If $X',Y'$ are copies of $X,Y$ respectively then $d[X';Y']=d[X ;Y]$.
\end{lemma}

\begin{proof} \uses{copy-ent}\leanok Immediate from Definitions \ref{ruz-dist-def} and \Cref{copy-ent}.
\end{proof}

\begin{lemma}[Ruzsa distance in independent case]\label{ruz-indep}
  \uses{ruz-dist-def}
  \lean{ProbabilityTheory.IndepFun.rdist_eq}\leanok
  If $X,Y$ are independent $G$-random variables then
  $$ d[X ;Y] := \bbH[X - Y] - \bbH[X]/2 - \bbH[Y]/2.$$
\end{lemma}

\begin{proof} \uses{relabeled-entropy, copy-ent}\leanok Immediate from \Cref{ruz-dist-def} and Lemmas \ref{relabeled-entropy}, \ref{copy-ent}.
\end{proof}


\begin{lemma}[Distance controls entropy difference]\label{ruzsa-diff}
  \uses{ruz-dist-def}
  \lean{diff_ent_le_rdist}\leanok
  If $X,Y$ are $G$-valued random variables, then
$$|\bbH[X]-H[Y]| \leq 2 d[X ;Y].$$
\end{lemma}

\begin{proof} \uses{sumset-lower, neg-ent} \leanok Immediate from \Cref{sumset-lower} and \Cref{ruz-dist-def}, and also \Cref{neg-ent}.
\end{proof}


\begin{lemma}[Distance nonnegative]\label{ruzsa-nonneg}
  \lean{rdist_nonneg}\leanok
  If $X,Y$ are $G$-valued random variables, then
  $$ d[X ;Y] \geq 0.$$
\end{lemma}

\begin{proof}
  \uses{ruzsa-diff}
  \leanok
  Immediate from \Cref{ruzsa-diff}.
\end{proof}


\begin{lemma}[Projection entropy and distance]\label{dist-projection}\lean{ent_of_proj_le}\leanok
If $G$ is an additive group and $X$ is a $G$-valued random variable and $H\leq G$ is a finite subgroup then, with $\pi:G\to G/H$ the natural homomorphism we have (where $U_H$ is uniform on $H$)
\[\mathbb{H}(\pi(X))\leq 2d[X;U_H].\]
\end{lemma}
\begin{proof}
\uses{independent-exist, ruzsa-diff, chain-rule, shear-ent, submodularity, jensen-bound}\leanok
WLOG, we make $X$, $U_H$ independent (\Cref{independent-exist}).
Now by Lemmas \ref{submodularity}, \ref{shear-ent}, \ref{jensen-bound}
\begin{align*}
&\mathbb{H}(X-U_H|\pi(X)) \geq \mathbb{H}(X-U_H|X) &= \mathbb{H}(U_H) \\
&\mathbb{H}(X-U_H|\pi(X)) \leq \log |H| &= \mathbb{H}(U_H)
\end{align*}
By \Cref{chain-rule}
\[\mathbb{H}(X-U_H)=\mathbb{H}(\pi(X))+\mathbb{H}(X-U_H|\pi(X))=\mathbb{H}(\pi(X))+\mathbb{H}(U_H)\]
and therefore
\[d[X;U_H]=\mathbb{H}(\pi(X))+\frac{\mathbb{H}(U_H)-\mathbb{H}(X)}{2}.\]
Furthermore by \Cref{ruzsa-diff}
\[d[X;U_H]\geq \frac{\lvert \mathbb{H}(X)-\mathbb{H}(U_H)\rvert}{2}.\]
Adding these inequalities gives the result.
\end{proof}

\begin{lemma}[Improved Ruzsa triangle inequality]
  \label{ruzsa-triangle-improved}
  \lean{ent_of_diff_le}
  \leanok
  If $X,Y,Z$ are $G$-valued random variables on $\Omega$ with $(X,Y)$ independent of $Z$, then
  \begin{equation}\label{submod-explicit} \bbH[X - Y] \leq \bbH[X-Z] + \bbH[Z-Y] - \bbH[Z]\end{equation}
\end{lemma}

This is an improvement over the usual Ruzsa triangle inequality because $X,Y$ are not assumed to be independent.  However we will not utilize this improvement here.

\begin{proof}
  \uses{alt-submodularity,subadditive, relabeled-entropy, add-entropy}
  \leanok
  Apply \Cref{alt-submodularity} to obtain
  \[\bbH[X - Z, X - Y] + \bbH[Y, X - Y] \geq \bbH[X - Z, Y, X - Y] + \bbH[X - Y].\]
  Using
  \[\bbH[X - Z, X - Y] \leq \bbH[X - Z] + \bbH[Y - Z]\]
  (from \Cref{relabeled-entropy}, \Cref{subadditive}),
  \[\bbH[Y, X - Y] = \bbH[X, Y] \]
  (from \Cref{relabeled-entropy}), and
  \[\bbH[X - Z, Y, X - Y] = \bbH[X, Y, Z] = \bbH[X, Y] + \bbH[Z]\]
  (from \Cref{relabeled-entropy} and \Cref{add-entropy}) and rearranging, we indeed obtain~\eqref{submod-explicit}.
\end{proof}

\begin{lemma}[Ruzsa triangle inequality]
  \label{ruzsa-triangle}
  \lean{rdist_triangle}
  \leanok
  If $X,Y,Z$ are $G$-valued random variables, then
$$ d[X ;Y] \leq d[X ;Z] + d[Z;Y].$$
\end{lemma}

\begin{proof}\uses{ruz-copy, ruz-indep, independent-exist, ruzsa-triangle-improved}\leanok By \Cref{ruz-copy} and Lemmas \ref{independent-exist}, \ref{ruz-indep}, it suffices to prove this inequality assuming that $X,Y,Z$ are defined on the same space and are independent.  But then the claim follows from \Cref{ruzsa-triangle-improved} and \Cref{ruz-dist-def}.
\end{proof}

\begin{definition}[Conditioned Ruzsa distance]\label{cond-dist-def}
  \uses{ruz-dist-def}
  \lean{condRuzsaDist}\leanok
If $(X, Z)$ and $(Y, W)$ are random variables (where $X$ and $Y$ are $G$-valued) we define
$$ d[X  | Z; Y | W] := \sum_{z,w}  \bbP[Z=z]  \bbP[W=w] d[(X|Z=z); (Y|(W=w))].$$
similarly
$$ d[X ; Y | W] := \sum_{w}  \bbP[W=w] d[X ; (Y|(W=w))].$$
\end{definition}

\begin{lemma}[Alternate form of distance]\label{cond-dist-alt}
\lean{condRuzsaDist_of_copy, condRuzsaDist'_of_copy, condRuzsaDist_of_indep, condRuzsaDist'_of_indep}\leanok
  The expression $d[X  | Z;Y | W]$ is unchanged if $(X,Z)$ or $(Y,W)$ is replaced by a copy.  Furthermore, if $(X,Z)$ and $(Y,W)$ are independent, then
$$  d[X  | Z;Y | W] = \bbH[X-Y|Z,W] - \bbH[X|Z]/2 - \bbH[Y|W]/2$$
and similarly
$$  d[X ;Y | W] = \bbH[X-Y|W] - \bbH[X]/2 - \bbH[Y|W]/2.$$
\end{lemma}

\begin{proof}\uses{copy-ent, ruz-copy, ruz-indep, cond-dist-def, conditional-entropy-def}\leanok Straightforward thanks to \Cref{copy-ent}, \Cref{ruz-copy}, \Cref{ruz-indep}, \Cref{cond-dist-def}, \Cref{conditional-entropy-def}.
\end{proof}


\begin{theorem}[Improved entropy version of PFR]\label{entropy-pfr-improv}\lean{entropic_PFR_conjecture_improv}\leanok
  Let $G = \F_2^n$, and suppose that $X^0_1, X^0_2$ are $G$-valued random variables.
  Then there is some subgroup $H \leq G$ such that
  \[
    d[X^0_1;U_H] + d[X^0_2;U_H] \le 10 d[X^0_1;X^0_2],
  \]
  where $U_H$ is uniformly distributed on $H$.
  Furthermore, both $d[X^0_1;U_H]$ and $d[X^0_2;U_H]$ are at most $\tfrac{11}{2} d[X^0_1;X^0_2]$.
\end{theorem}

\chapter{Weak PFR over the integers}

\begin{lemma}\label{torsion-free-doubling}\lean{torsion_free_doubling}\leanok
If $G$ is torsion-free and $X,Y$ are $G$-valued random variables then $d[X;2Y]\leq 5d[X;Y]$.
\end{lemma}
\begin{proof}
\uses{alt-submodularity,ruzsa-triangle,ruzsa-diff}\leanok
Let $Y_1,Y_2$ be independent copies of $Y$ (also independent of $X$). Since $G$ is torsion-free we know $X,Y_1-Y_2,X-2Y_1$ uniquely determine $X,Y_1,Y_2$ and so
\[\mathbb{H}(X,Y_1,Y_2,X-2Y_1)=\mathbb{H}(X,Y_1,Y_2)=\mathbb{H}(X)+2\mathbb{H}(Y).\]
Similarly
\[\mathbb{H}(X,X-2Y_1)=\mathbb{H}(X)+\mathbb{H}(2Y_1)=\mathbb{H}(X)+\mathbb{H}(Y).\]
Furthermore
\[\mathbb{H}(Y_1-Y_2,X-2Y_1)=\mathbb{H}(Y_1-Y_2,X-Y_1-Y_2)\leq \mathbb{H}(Y_1-Y_2)+\mathbb{H}(X-Y_1-Y_2).\]
By submodularity (\Cref{alt-submodularity})
\[\mathbb{H}(X,Y_1,Y_2,X-2Y_1)+\mathbb{H}(X-2Y_1)\leq \mathbb{H}(X,X-2Y_1)+\mathbb{H}(Y_1-Y_2,X-2Y_1).\]
Combining these inequalities
\[\mathbb{H}(X-2Y_1)\leq \mathbb{H}(Y_1-Y_2)+\mathbb{H}(X-Y_1-Y_2)-\mathbb{H}(Y).\]
Similarly we have
\[\mathbb{H}(Y_1,Y_2,X-Y_1-Y_2)=\mathbb{H}(X)+2\mathbb{H}(Y),\]
\[\mathbb{H}(Y_1,X-Y_1-Y_2)=\mathbb{H}(Y)+\mathbb{H}(X-Y_2),\]
and
\[\mathbb{H}(Y_2,X-Y_1-Y_2)=\mathbb{H}(Y)+\mathbb{H}(X-Y_1)\]
and by submodularity (\Cref{alt-submodularity}) again
\[\mathbb{H}(Y_1,Y_2,X-Y_1-Y_2)+ \mathbb{H}(X-Y_1-Y_2)\leq \mathbb{H}(Y_1,X-Y_1-Y_2)+\mathbb{H}(Y_2,X-Y_1-Y_2).\]
Combining these inequalities (and recalling the definition of Ruzsa distance) gives
\[\mathbb{H}(X-Y_1-Y_2)\leq \mathbb{H}(X-Y_1)+\mathbb{H}(X-Y_2)-\mathbb{H}(X)=2d[X;Y]+\mathbb{H}(Y).\]
It follows that
\[\mathbb{H}(X-2Y_1)\leq \mathbb{H}(Y_1-Y_2)+2d[X;Y]\]
and so (using $\mathbb{H}(2Y)=\mathbb{H}(Y)$)
\begin{align*}
d[X;2Y]
&=\mathbb{H}(X-2Y_1)-\mathbb{H}(X)/2-\mathbb{H}(2Y)/2\\
&\leq \mathbb{H}(Y_1-Y_2)+2d[X;Y]-\mathbb{H}(X)/2-\mathbb{H}(Y)/2\\
&= d[Y_1;Y_2]+\frac{\mathbb{H}(Y)-\mathbb{H}(X)}{2}+2d[X;Y].
\end{align*}
Finally note that by the triangle inequality (\Cref{ruzsa-triangle}) we have
\[d[Y_1;Y_2]\leq d[Y_1;X]+d[X;Y_2]=2d[X;Y].\]
The result follows from $(\mathbb{H}(Y)-\mathbb{H}(X))/2\leq d[X;Y]$ (\Cref{ruzsa-diff}).
\end{proof}

\begin{lemma}\label{torsion-dist-shrinking}\lean{torsion_dist_shrinking}\leanok
If $G$ is a torsion-free group and $X,Y$ are $G$-valued random variables and $\phi:G\to \mathbb{F}_2^d$ is a homomorphism then
\[\mathbb{H}(\phi(X))\leq 10d[X;Y].\]
\end{lemma}
\begin{proof}
\uses{fibring-ineq, torsion-free-doubling, dist-zero}\leanok
By \Cref{fibring-ineq} and \Cref{torsion-free-doubling} we have
\[d[\phi(X);\phi(2Y)]\leq d[X;2Y]\leq 5d[X;Y]\]
and $\phi(2Y)=2\phi(Y)\equiv 0$ so the left-hand side is equal to $d[\phi(X);0]=\mathbb{H}(\phi(X))/2$ (using \Cref{dist-zero}).
\end{proof}

\begin{lemma}\label{app-ent-pfr}\lean{app_ent_PFR}\leanok
Let $G=\mathbb{F}_2^n$ and $\alpha\in (0,1)$ and let $X,Y$ be $G$-valued
random variables such that
\[\mathbb{H}(X)+\mathbb{H}(Y)> \frac{20}{\alpha} d[X;Y].\]
There is a non-trivial subgroup $H\leq G$ such that
\[\log \lvert H\rvert <\frac{1+\alpha}{2}(\mathbb{H}(X)+\mathbb{H}(Y))\] and
\[\mathbb{H}(\psi(X))+\mathbb{H}(\psi(Y))< \alpha (\mathbb{H}(X)+\mathbb{H}(Y))\]
where $\psi:G\to G/H$ is the natural projection homomorphism.
\end{lemma}
\begin{proof}
\uses{entropy-pfr-improv, ruzsa-diff, dist-projection, ruzsa-nonneg}\leanok
By \Cref{entropy-pfr-improv} there exists a subgroup $H$ such that
$d[X;U_H] + d[Y;U_H] \leq 10 d[X;Y]$. Using \Cref{dist-projection} we
deduce that $\mathbb{H}(\psi(X)) + \mathbb{H}(\psi(X)) \leq 20 d[X;Y]$. The
second claim follows adding these inequalities and using the assumption on
$\mathbb{H}(X)+\mathbb{H}(Y)$.

Furthermore we have by \Cref{ruzsa-diff}
\[\log \lvert H \rvert-\mathbb{H}(X)\leq 2d[X;U_H]\]
and similarly for $Y$ and thus
\begin{align*}
\log \lvert  H\rvert
&\leq
\frac{\mathbb{H}(X)+\mathbb{H}(Y)}{2}+d[X;U_H] + d[Y;U_H] \leq
\frac{\mathbb{H}(X)+\mathbb{H}(Y)}{2}+ 10d[X;Y]
\\& <
\frac{1+\alpha}{2}(\mathbb{H}(X)+\mathbb{H}(Y)).
\end{align*}
Finally note that if $H$
were trivial then $\psi(X)=X$ and $\psi(Y)=Y$ and hence
$\mathbb{H}(X)+\mathbb{H}(Y)=0$, which contradicts \Cref{ruzsa-nonneg}.
\end{proof}


\begin{lemma}\label{pfr-projection'}\lean{PFR_projection'}\leanok
If $G=\mathbb{F}_2^d$ and $\alpha\in (0,1)$ and $X,Y$ are $G$-valued random
variables then there is a subgroup $H\leq \mathbb{F}_2^d$ such that
\[\log \lvert H\rvert \leq \frac{1+\alpha}{2(1-\alpha)} (\mathbb{H}(X)+\mathbb{H}(Y))\]
and if $\psi:G \to G/H$ is the natural projection then
\[\mathbb{H}(\psi(X))+\mathbb{H}(\psi(Y))\leq \frac{20}{\alpha} d[\psi(X);\psi(Y)].\]
\end{lemma}
\begin{proof}
\uses{app-ent-pfr}\leanok
Let $H\leq \mathbb{F}_2^d$ be a maximal subgroup such that
\[\mathbb{H}(\psi(X))+\mathbb{H}(\psi(Y))> \frac{20}{\alpha} d[\psi(X);\psi(Y)]\]
and such that there exists $c \ge 0$ with
\[\log \lvert H\rvert \leq \frac{1+\alpha}{2(1-\alpha)}(1-c)(\mathbb{H}(X)+\mathbb{H}(Y))\]
and
\[\mathbb{H}(\psi(X))+\mathbb{H}(\psi(Y))\leq c (\mathbb{H}(X)+\mathbb{H}(Y)).\]
Note that this exists since $H=\{0\}$ is an example of such a subgroup or we are done with this choice of $H$.

We know that $G/H$ is a $2$-elementary group and so by Lemma
\ref{app-ent-pfr} there exists some non-trivial subgroup $H'\leq G/H$ such
that
\[\log \lvert H'\rvert < \frac{1+\alpha}{2}(\mathbb{H}(\psi(X))+\mathbb{H}(\psi(Y))\]
and
\[\mathbb{H}(\psi' \circ\psi(X))+\mathbb{H}(\psi' \circ \psi(Y))< \alpha(\mathbb{H}(\psi(X))+\mathbb{H}(\psi(Y)))\]
where $\psi':G/H\to (G/H)/H'$. By group isomorphism theorems we know that
there exists some $H''$ with $H\leq H''\leq G$ such that $H'\cong H''/H$ and
$\psi' \circ \psi(X)=\psi''(X)$ where $\psi'':G\to G/H''$ is the projection
homomorphism.

Since $H'$ is non-trivial we know that $H$ is a proper subgroup of $H''$. On the other hand we know that
\[\log \lvert H''\rvert=\log \lvert H'\rvert+\log \lvert H\rvert< \frac{1+\alpha}{2(1-\alpha)}(1-\alpha c)(\mathbb{H}(X)+\mathbb{H}(Y))\]
and
\[\mathbb{H}(\psi''(X))+\mathbb{H}(\psi''(Y))< \alpha (\mathbb{H}(\psi(X))+\mathbb{H}(\psi(Y)))\leq \alpha c (\mathbb{H}(X)+\mathbb{H}(Y)).\]
Therefore (using the maximality of $H$) it must be the first condition that fails, whence
\[\mathbb{H}(\psi''(X))+\mathbb{H}(\psi''(Y))\leq \frac{20}{\alpha}d[\psi''(X);\psi''(Y)].\]
\end{proof}

We could use the previous lemma for any value of $\alpha \in (0,1)$, which
would give a whole range of estimates in Theorem~\ref{weak-pfr-int}. For
definiteness, we specialize only to $\alpha=3/5$, which gives a constant $2$
in the first bound below.

\begin{lemma}\label{pfr-projection}\lean{PFR_projection}\leanok
If $G=\mathbb{F}_2^d$ and $\alpha\in (0,1)$ and $X,Y$ are $G$-valued random
variables then there is a subgroup $H\leq \mathbb{F}_2^d$ such that
\[\log \lvert H\rvert \leq 2 (\mathbb{H}(X)+\mathbb{H}(Y))\]
and if $\psi:G \to G/H$ is the natural projection then
\[\mathbb{H}(\psi(X))+\mathbb{H}(\psi(Y))\leq 34 d[\psi(X);\psi(Y)].\]
\end{lemma}
\begin{proof}
\uses{pfr-projection'}\leanok
Specialize \Cref{pfr-projection'} to $\alpha=3/5$. In the second
inequality, it gives a bound $100/3 < 34$.
\end{proof}



\begin{lemma}\label{single-fibres}\lean{single_fibres}\leanok
Let $\phi:G\to H$ be a homomorphism and $A,B\subseteq G$ be finite subsets. If $x,y\in H$ then let $A_x=A\cap \phi^{-1}(x)$ and $B_y=B\cap \phi^{-1}(y)$. There exist $x,y\in H$ such that $A_x,B_y$ are both non-empty and
\[d[\phi(U_A);\phi(U_B)]\log \frac{\lvert A\rvert\lvert B\rvert}{\lvert A_x\rvert\lvert B_y\rvert}\leq (\mathbb{H}(\phi(U_A))+\mathbb{H}(\phi(U_B)))(d(U_A,U_B)-d(U_{A_x},U_{B_y})).\]
\end{lemma}
\begin{proof}
\uses{fibring-ident}\leanok
The random variables $(U_A\mid \phi(U_A)=x)$ and $(U_B\mid \phi(U_B)=y)$ are equal in distribution to $U_{A_x}$ and $U_{B_y}$ respectively (both are uniformly distributed over their respective fibres). It follows from \Cref{fibring-ident} that
\begin{align*}
\sum_{x,y\in H}\frac{\lvert A_x\rvert\lvert B_y\rvert}{\lvert A\rvert\lvert B\rvert}d[U_{A_x};U_{B_y}]
&=d[U_A\mid \phi(U_A); U_B\mid \phi(U_B)]\\
&\leq d[U_A;U_B]-d[\phi(U_A);\phi(U_B)].
\end{align*}
Therefore with $M:=\mathbb{H}(\phi(U_A))+\mathbb{H}(\phi(U_B))$ we have
\[\left(\sum_{x,y\in H}\frac{\lvert A_x\rvert\lvert B_y\rvert}{\lvert A\rvert\lvert B\rvert}Md[U_{A_x};U_{B_y}]\right)+Md[\phi(U_A);\phi(U_B)]\leq Md[U_A;U_B].\]
Since
\[M=\sum_{x,y\in H}\frac{\lvert A_x\rvert\lvert B_y\rvert}{\lvert A\rvert\lvert B\rvert}\log \frac{\lvert A\rvert\lvert B\rvert}{\lvert A_x\rvert\lvert B_y\rvert}\]
we have
\[\sum_{x,y\in H} \frac{\lvert A_x\rvert\lvert B_y\rvert}{\lvert A\rvert\lvert B\rvert}\left(Md[U_{A_x};U_{B_y}]+d[\phi(U_A);\phi(U_B)]\log \frac{\lvert A\rvert\lvert B\rvert}{\lvert A_x\rvert\lvert B_y\rvert}\right)\leq  Md[U_A;U_B].\]
It follows that there exists some $x,y\in H$ such that $\lvert A_x\rvert,\lvert B_y\rvert\neq 0$ and
\[Md[U_{A_x};U_{B_y}]+d[\phi(U_A);\phi(U_B)]\log \frac{\lvert A\rvert\lvert B\rvert}{\lvert A_x\rvert\lvert B_y\rvert}\leq  Md[U_A;U_B].\]
\end{proof}




\begin{theorem}\label{weak-pfr-int}\lean{weak_PFR_int}\leanok
Let $A\subseteq \mathbb{Z}^d$ and $\lvert A-A\rvert\leq K\lvert A\rvert$.
There exists $A'\subseteq A$ such that $\lvert A'\rvert \geq K^{-17}\lvert
A\rvert$ and $\dim A' \leq \frac{40}{\log 2}\log K$.
\end{theorem}

\begin{theorem}\label{weak-pfr-asymm}\lean{weak_PFR_asymm}\leanok
If $A,B\subseteq \mathbb{Z}^d$ are finite non-empty sets then there exist non-empty $A'\subseteq A$ and $B'\subseteq B$ such that
\[\log\frac{\lvert A\rvert\lvert B\rvert}{\lvert A'\rvert\lvert B'\rvert}\leq 34d[U_A;U_B]\]
such that $\max(\dim A',\dim B')\leq \frac{40}{\log 2} d[U_A;U_B]$.
\end{theorem}

\begin{proof}
\uses{torsion-dist-shrinking, pfr-projection, single-fibres, dimension-def}\leanok
Without loss of generality we can assume that $A$ and $B$ are not both inside (possibly distinct) cosets of the same subgroup of $\mathbb{Z}^d$, or we just replace $\mathbb{Z}^d$ with that subgroup. We prove the result by induction on $\lvert A\rvert+\lvert B\rvert$.

Let $\phi:\mathbb{Z}^d\to \mathbb{F}_2^d$ be the natural mod-2 homomorphism. By \Cref{torsion-dist-shrinking}
\[\max(\mathbb{H}(\phi(U_A)),\mathbb{H}(\phi(U_B)))\leq 10d[U_A;U_B].\]
We now apply \Cref{pfr-projection}, obtaining some subgroup $H\leq \mathbb{F}_2^d$ such that
\[\log \lvert H\rvert \leq 40d[U_A;U_B]\]
and
\[\mathbb{H}(\tilde{\phi}(U_A))+\mathbb{H}(\tilde{\phi}(U_B))\leq 34 d[\tilde{\phi}(U_A);\tilde{\phi}(U_B)]\]
where $\tilde{\phi}:\mathbb{Z}^d\to \mathbb{F}_2^d/H$ is $\phi$ composed with the projection onto $\mathbb{F}_2^d/H$.


By \Cref{single-fibres} there exist $x,y\in \mathbb{F}_2^d/H$ such that, with $A_x=A\cap \tilde{\phi}^{-1}(x)$ and similarly for $B_y$,
\[\log \frac{\lvert A\rvert\lvert B\rvert}{\lvert A_x\rvert\lvert B_y\rvert}\leq 34(d[U_A;U_B]-d[U_{A_x};U_{B_y}]).\]
Suppose first that $\lvert A_x\rvert+\lvert B_y\rvert=\lvert A\rvert+\lvert B\rvert$. This means that $\tilde{\phi}(A)=\{x\}$ and $\tilde{\phi}(B)=\{y\}$, and hence both $A$ and $B$ are in cosets of $\ker \tilde{\phi}$. Since by assumption $A,B$ are not in cosets of a proper subgroup of $\mathbb{Z}^d$ this means $\ker\tilde{\phi}=\mathbb{Z}^d$, and so (examining the definition of $\tilde{\phi}$) we must have $H=\mathbb{F}_2^d$. Then our bound on $\log\lvert H\rvert$ forces $d\leq \frac{40}{\log 2}d[U_A;U_B]$ and we are done with $A'=A$ and $B'=B$.

Otherwise,
\[\lvert A_x\rvert+\lvert B_y\rvert <\lvert A\rvert+\lvert B\rvert.\]
By induction we can find some $A'\subseteq A_x$ and $B'\subseteq B_y$ such that $\dim A',\dim B'\leq \frac{40}{\log 2} d[U_{A_x};U_{B_y}]\leq \frac{40}{\log 2}d[U_A;U_B]$ and

\[\log \frac{\lvert A_x\rvert\lvert B_y\rvert}{\lvert A'\rvert\lvert B'\rvert}\leq 34d[U_{A_x};U_{B_y}].\]
Adding these inequalities implies

\[\log\frac{\lvert A\rvert\lvert B\rvert}{\lvert A'\rvert\lvert B'\rvert}\leq 34d[U_A;U_B]\]
as required.
\end{proof}

\end{document}
